\documentclass[17pt]{beamer}
%\documentclass{beamer}
\usetheme{Madrid}
\useoutertheme{noslidenum}
\setbeamertemplate{navigation symbols}{}
\usepackage{beamerthemesplit}
\usepackage{ae,aecompl}
\usepackage[scaled=.95]{helvet}
\usepackage[english]{babel}
\usepackage[latin1]{inputenc}
\usepackage[T1]{fontenc}
\usepackage{mathpazo,courier,euler}
\usepackage{listings}
\lstset{language=sh,
    basicstyle=\ttfamily\bfseries,
  showstringspaces=false,
  keywordstyle=\color{black}\bfseries}

% theme split
\usepackage{verbatim}
\newenvironment{colorverbatim}[1][]%
{%
\color{blue}
\verbatim
}%
{%
\endverbatim
}%

\definecolor{NavyBlue}{RGB}{0, 76, 153}
\setbeamercolor{structure}{fg=NavyBlue}
\author[FOSSEE]{}
\institute[IIT Bombay]{}
\date[]{}

% logo
\logo{\includegraphics[height=1.30cm]{../images/3t-logo.pdf}}
\logo{\includegraphics[height=0.6cm]{../images/fossee-logo.png}

\hspace{200pt}
\includegraphics[scale=0.5]{../images/fossee-logo.png}\\
\hspace{275pt}
\includegraphics[scale=0.9]{../images/3t-logo.pdf}}

\begin{document}

\sffamily \bfseries
\title
[Getting Started With Lists]
{Getting Started With Lists}
\author
[FOSSEE, IIT - Bombay]
{\small Talk to a Teacher\\{\color{blue}\url{http://spoken-tutorial.org}}\\National Mission on Education
 through ICT\\{\color{blue}\url{http://sakshat.ac.in}} \\ [0.8cm]Script by: Hardik Ghaghada \\ Narration by: \\
 [0.5cm]
{\small 15 May 2013}}

%%%%%%%%%%%%%%%%%%%%%%%%%%%%%%%%%%%%%%%%%%%%%%%%%%%%%%%%%%%%%%%%%%%%%%%
\begin{frame}
  \titlepage
\end{frame}
%%%%%%%%%%%%%%%%%%%%%%%%%%%%%%%%%%%%%%%%%%%%%%%%%%%%%%%%%%%%%%%%%%%%%%%
\begin{frame}
\frametitle{Objectives}
  At the end of this tutorial, you will be able to, 
\begin{itemize}
\item Create lists
\item Access list elements
\item Append elements to lists
\item Delete elements from lists
\end{itemize}
\end{frame}
%%%%%%%%%%%%%%%%%%%%%%%%%%%%%%%%%%%%%%%%%%%%%%%%%%%%%%%%%%%%%%%%%%%%%%%
\begin{frame}
\frametitle{Lists}
List is a compound data type, it can contain data of mutually
different datatypes. List is also a sequence data type where all the
elements are arranged in a specific order.
\end{frame}
%%%%%%%%%%%%%%%%%%%%%%%%%%%%%%%%%%%%%%%%%%%%%%%%%%%%%%%%%%%%%%%%%%%%%%%
\begin{frame}
\frametitle{Exercise 1}
\begin{itemize}
\item What happens when you do nonempty[-1].
\end{itemize}
\end{frame}
%%%%%%%%%%%%%%%%%%%%%%%%%%%%%%%%%%%%%%%%%%%%%%%%%%%%%%%%%%%%%%%%%%%%%%%
\begin{frame}
\frametitle{Exercise 2}
\begin{itemize}
\item What is the syntax to get the element `and' 
     in the list,listinlist ?
\item How would you get `and' using negative indices?
\end{itemize}
\end{frame}
%%%%%%%%%%%%%%%%%%%%%%%%%%%%%%%%%%%%%%%%%%%%%%%%%%%%%%%%%%%%%%%%%%%%%%%
\begin{frame}[fragile]
\frametitle{Solution 2}
\lstset{language=Python}
\begin{lstlisting}
listinlist[1]
listinlist[-5]
\end{lstlisting}
\end{frame}
%%%%%%%%%%%%%%%%%%%%%%%%%%%%%%%%%%%%%%%%%%%%%%%%%%%%%%%%%%%%%%%%%%%%%%%
\begin{frame}
\frametitle{Exercise 3}
\begin{itemize}
\item Remove the third element from the list, listinlist.
\item Remove `and' from the list, listinlist.
\end{itemize}
\end{frame}
%%%%%%%%%%%%%%%%%%%%%%%%%%%%%%%%%%%%%%%%%%%%%%%%%%%%%%%%%%%%%%%%%%%%%%%
\begin{frame}[fragile]
\frametitle{Solution 3}
\lstset{language=Python}
\begin{lstlisting}
del(listinlist[2])
listinlist.remove('and')
\end{lstlisting}
\end{frame}
%%%%%%%%%%%%%%%%%%%%%%%%%%%%%%%%%%%%%%%%%%%%%%%%%%%%%%%%%%%%%%%%%%%%%%%
\begin{frame}[fragile]
	\frametitle{Summary}
	In this tutorial we have learnt how to,
    \begin{itemize}
        \item Create lists
        \item Access elements
        \item Use ``append'' function
        \item Use ``del'' function
        \item Use ``remove'' function.
        \item Use ``len'' function.
	\end{itemize}
\end{frame}
%%%%%%%%%%%%%%%%%%%%%%%%%%%%%%%%%%%%%%%%%%%%%%%%%%%%%%%%%%%%%%%%%%%%%%%
\begin{frame}[fragile]
	\frametitle{Evaluation}
	\begin{enumerate}
        \item How do you create an empty list ?
        \vspace{8pt}
        \item Can you have a list inside a list ?
        \vspace{8pt}
        \item How would you access the last element of a list without finding its length?
    \end{enumerate}
\end{frame}
%%%%%%%%%%%%%%%%%%%%%%%%%%%%%%%%%%%%%%%%%%%%%%%%%%%%%%%%%%%%%%%%%%%%%%%
\begin{frame}[fragile]
	\frametitle{Solutions}
	\begin{enumerate}
        \item empty=[]
        \vspace{8pt}
        \item Yes we can have list inside list\\
        \vspace{8pt}
        \item Use negative index to access the last element of the list
    \end{enumerate}
\end{frame}
%%%%%%%%%%%%%%%%%%%%%%%%%%%%%%%%%%%%%%%%%%%%%%%%%%%%%%%%%%%%%%%%%%%%%%%
\begin{frame}
\frametitle{FOSSEE}
{\color{NavyBlue}Free and Open Source Software for \\Education} \\
\begin{itemize}
\item To enable the use of free and open source software tools
\item For more details, please visit: {\color{blue}\url{http://fossee.in/}}
\end{itemize}
\end{frame}
%%%%%%%%%%%%%%%%%%%%%%%%%%%%%%%%%%%%%%%%%%%%%%%%%%%%%%%%%%%%%%%%%%%%%%%
\begin{frame}
\frametitle{About the Spoken Tutorial Project}
\begin{itemize}
\item Watch the video available at {\color{blue}\url{http://spoken-tutorial.org /What\_is\_a\_Spoken\_Tutorial}} 
\item It summarises the Spoken Tutorial project 
\item If you do not have good bandwidth, you can download and watch it
\end{itemize}
\end{frame}
%%%%%%%%%%%%%%%%%%%%%%%%%%%%%%%%%%%%%%%%%%%%%%%%%%%%%%%%%%%%%%%%%%%%%%%
\begin{frame}
\frametitle{Spoken Tutorial Workshops}The Spoken Tutorial Project Team 
\begin{itemize}
\item Conducts workshops using spoken tutorials 
\item Gives certificates to those who pass an online test 
\item For more details, please write to \\ \hspace {0.5cm}{\color{blue}contact@spoken-tutorial.org}
\end{itemize}
\end{frame}
%%%%%%%%%%%%%%%%%%%%%%%%%%%%%%%%%%%%%%%%%%%%%%%%%%%%%%%%%%%%%%%%%%%%%%%
\begin{frame}
\frametitle{Acknowledgements}
\begin{itemize}
\item Spoken Tutorial Project is a part of the Talk to a Teacher  project 
\item Supported by the National Mission on Education through  ICT, MHRD, Government of India 
\item More information: \\{\color{blue}\url{http://spoken-tutorial.org/NMEICT-Intro}}
\end{itemize}
\end{frame}
%%%%%%%%%%%%%%%%%%%%%%%%%%%%%%%%%%%%%%%%%%%%%%%%%%%%%%%%%%%%%%%%%%%%%%%
\begin{frame}

  \begin{block}{}
  \begin{center}
  \textcolor{blue}{\Large THANK YOU!} 
  \end{center}
  \end{block}
\begin{block}{}
  \begin{center}
    For more Information, visit our website\\
    {\color{blue}\url{http://fossee.in/}}
  \end{center}  
  \end{block}
\end{frame}
%%%%%%%%%%%%%%%%%%%%%%%%%%%%%%%%%%%%%%%%%%%%%%%%%%%%%%%%%%%%%%%%%%%%%%%
\end{document}
